\section{Introduction} % possibile estensione sullo stato dell'arte
Self-Protecting Systems (SPS) are autonomic systems that detect and mitigate
cyber threats with limited or no human intervention~\cite{yuan2014systematic}.
In the canonical SPS architecture (see Figure~\ref{fig:architecture}), a managed
system is supervised by an autonomic manager responsible for monitoring,
detection, and response. Although effective in many settings, this design also
introduces a structural vulnerability: an adversary may target the
autonomic manager itself, compromising its ability to detect attacks and
coordinate responses~\cite{chess2003security}. In SPS research, protecting the
component that protects the system remains a central and still largely open
challenge. Prior work~\cite{11126595} addressed this weakness by implementing
the autonomic manager as a permissioned blockchain. In that approach, alerts
generated by Intrusion Detection Systems (IDS) are submitted as transactions to
a blockchain operated by multiple system nodes and recorded in the state of a
smart contract. The underlying consensus mechanism provides a consistent,
system-wide view of the monitored state. The same contract also determines
which mitigation actions should be executed and when. Actuator nodes then
retrieve these decisions from the blockchain and enforce the corresponding
actions.
The key idea for improving the protection of the autonomic manager is that
consensus tolerates a bounded number of compromised replicas, therefore, the
manager can continue to operate correctly even under partial compromise.
Moreover, the smart-contract logic can be designed to mitigate malicious inputs
(e.g., by requiring corroboration from multiple independent alerts before
triggering a response).

\begin{figure}[t!]
  \centering
  \includegraphics[width=0.8\textwidth]{autonomic_arch.png}
  \caption{Canonical Self-Protecting System architecture. In this architecture, blockchain-related technologies may be adopted as part of the autonomic manager.}
  \label{fig:architecture}
\end{figure}

We stress that this is a rather atypical use of blockchain technology, and we argue
that, in this specific setting, selecting the underlying blockchain platform
is not a neutral choice.
In~\cite{11126595}, Quorum is adopted as a natural baseline choice.
However, despite being a well-established technology for permissioned settings, it is not automatically
a good choice in our specific context.

In this work, we revisit that choice by analyzing the specific needs of the SPS manager and comparing them with the features of Quorum and CometBFT~\cite{cometbft}.  
CometBFT is a lower-level technology designed to provide a highly customizable Distributed State Machine (DSM).

The contribution of this paper is threefold:  
(i)~we identify a set of requirements for applying blockchain technology to SPS systems,  
(ii)~we compare Quorum and CometBFT against those requirements, and  
(iii)~we present initial experimental comparison of Quorum vs. CometBFT.

The remainder of this paper is organized as follows.  
Section~\ref{sec:architecture} describes the SPS context and briefly presents the architecture in which we plan to adopt blockchain technology.  
Section~\ref{sec:requirements} defines the blockchain requirements in this specific application context.  
Section~\ref{sec:comparison} compares Quorum and CometBFT with respect to those requirements and discusses the advantages of CometBFT over Quorum.  
Section~\ref{sec:experiments} reports an initial experimental comparison of Quorum and CometBFT.
Finally, Section~\ref{sec:conclusions} concludes the paper.



% In this work we aim at analyzing this choice to make 

% We make three contributions: (i)~we identify the minimal distributed-systems properties required for cyber-resilient SPS coordination, (ii)~we show that certain typical blockchain features are unnecessary in our application setting and introduce unneeded complexity, and (iii)~we aim at proposing a requirement-driven architecture based on a Distributed State Machine (DSM) for the coordination layer, which essentially takes only the needed part of a blockchain architecture. In this perspective, we further outline an evaluation of established BFT consensus engines (e.g., Tendermint~\cite{GitHubte7:online}) as a concrete instantiation of the DSM substrate in our SPS application setting.